% GENERATED FILE. Edit canonical lessons, then run npm run generate:books.
% canonical-source-hash: fnv1a64:ea8d47f7db206659
% canonical-lessons: SA-C393-vijayate, SA-C393-avicalah, SA-C393-nirantaram, SA-C393-tatparah, SA-C393-krtakrtyah, SA-R393-review-387-393, SA-R393-review-final-393

\chapter{What It Takes}
\label{ch:sa-what-it-takes}
\hlchaptermodality{393}

\begin{chapteropening}
\textbf{By the end of this chapter:} \emph{I can say what it will take to succeed.}

Complete the last lesson of chapter 393: I can say what it will take to succeed.
\end{chapteropening}

\section[\texorpdfstring{\sk{विजयते} (vijayate)}{विजयते (vijayate)}]{\sk{विजयते} (vijayate) — triumphs; wins through}

\label{lesson:SA-C393-vijayate}



\begin{quote}
Before the new one: say the Sanskrit for co-operation, then the Sanskrit for a kindness.
\end{quote}

\subsection*{You'll want to know: \sk{विजयते}}
\textbf{\sk{विजयते}} — \emph{vijayate} — \textquotedblleft{}triumphs; wins through\textquotedblright{}.

\begin{quote}
\sk{यदि} \sk{धैर्यम्} \sk{अस्ति}, \sk{तर्हि} \sk{सः} \sk{अवश्यं} \sk{विजयते।}
\end{quote}

\begin{quote}
\emph{yadi dhairyam asti, tarhi saḥ avaśyaṃ vijayate.}
\end{quote}

\begin{quote}
If he has courage, he will surely win through.
\end{quote}

Past (he …): \sk{सः} \sk{शत्रुं} \sk{विजितवान्}, \emph{saḥ śatruṃ vijitavān}.

\begin{cousinweb}
\textbf{\sk{विजयते}} (\emph{vijayate}) is built on \emph{vi} and the root √ji, \textquotedblleft{}to win\textquotedblright{}; MW gives \emph{vi-√ji} \textquotedblleft{}to conquer, to be victorious\textquotedblright{} (MW).
\end{cousinweb}

\begin{grammarlens}[title={What you've built}]
One more word for saying what happens if.
\end{grammarlens}

\subsection*{Your turn}
\begin{itemize}
\raggedright
  \item \emph{Say it:} \emph{vijayate}
  \item \emph{Say it:} \emph{vijayate}, once more
  \item \emph{Say it:} \emph{yadi dhairyam asti, tarhi saḥ avaśyaṃ vijayate.} — If he has courage, he will surely win through.
  \item \emph{Recall:} say the Sanskrit for co-operation, then the Sanskrit for a kindness, then say \emph{vijayate} again
\end{itemize}

\begin{tcolorbox}[breakable,colback=teal!4,colframe=teal!35!black,title={Before you move on}]
What does \sk{विजयते} mean? (\textquotedblleft{}Triumphs; wins through\textquotedblright{}.) Say it once more.
\end{tcolorbox}
\section[\texorpdfstring{\sk{अविचलः} (avicalaḥ)}{अविचलः (avicalaḥ)}]{\sk{अविचलः} (avicalaḥ) — unshakable; steady}

\label{lesson:SA-C393-avicalah}



\begin{quote}
Before the new one: say the Sanskrit for a kindness, then the Sanskrit for triumphs.
\end{quote}

\subsection*{You'll want to know: \sk{अविचलः}}
\textbf{\sk{अविचलः}} — \emph{avicalaḥ} — \textquotedblleft{}unshakable; steady\textquotedblright{}.

\begin{quote}
\sk{यदि} \sk{मनः} \sk{अविचलम्}, \sk{तर्हि} \sk{भयं} \sk{न} \sk{भवति।}
\end{quote}

\begin{quote}
\emph{yadi manaḥ avicalam, tarhi bhayaṃ na bhavati.}
\end{quote}

\begin{quote}
If the mind is unshakable, there is no fear.
\end{quote}

\begin{cousinweb}
\textbf{\sk{अविचलः}} (\emph{avicalaḥ}) is built from \emph{a-}, \textquotedblleft{}not\textquotedblright{}, and \emph{vi-cala}, \textquotedblleft{}moving about\textquotedblright{}, from the root √cal, \textquotedblleft{}to move\textquotedblright{}; MW gives \emph{a-vicala} \textquotedblleft{}immovable, steady, firm\textquotedblright{} (MW).
\end{cousinweb}

\begin{grammarlens}[title={What you've built}]
One more word for saying what happens if.
\end{grammarlens}

\subsection*{Your turn}
\begin{itemize}
\raggedright
  \item \emph{Say it:} \emph{avicalaḥ}
  \item \emph{Say it:} \emph{avicalaḥ}, once more
  \item \emph{Say it:} \emph{yadi manaḥ avicalam, tarhi bhayaṃ na bhavati.} — If the mind is unshakable, there is no fear.
  \item \emph{Recall:} say the Sanskrit for a kindness, then the Sanskrit for triumphs, then say \emph{avicalaḥ} again
\end{itemize}

\begin{tcolorbox}[breakable,colback=teal!4,colframe=teal!35!black,title={Before you move on}]
What does \sk{अविचलः} mean? (\textquotedblleft{}Unshakable; steady\textquotedblright{}.) Say it once more.
\end{tcolorbox}
\section[\texorpdfstring{\sk{निरन्तरम्} (nirantaram)}{निरन्तरम् (nirantaram)}]{\sk{निरन्तरम्} (nirantaram) — without a break; constantly}

\label{lesson:SA-C393-nirantaram}



\begin{quote}
Before the new one: say the Sanskrit for triumphs, then the Sanskrit for unshakable.
\end{quote}

\subsection*{You'll want to know: \sk{निरन्तरम्}}
\textbf{\sk{निरन्तरम्}} — \emph{nirantaram} — \textquotedblleft{}without a break; constantly\textquotedblright{}.

\begin{quote}
\sk{यदि} \sk{भवान्} \sk{निरन्तरं} \sk{पठति}, \sk{तर्हि} \sk{परीक्षायाम्} \sk{उत्तमः} \sk{भविष्यति।}
\end{quote}

\begin{quote}
\emph{yadi bhavān nirantaraṃ paṭhati, tarhi parīkṣāyām uttamaḥ bhaviṣyati.}
\end{quote}

\begin{quote}
If you study without a break, you will do well in the exam.
\end{quote}

\begin{cousinweb}
\textbf{\sk{निरन्तरम्}} (\emph{nirantaram}) is built from \emph{nir}, \textquotedblleft{}without\textquotedblright{}, and \emph{antara}, \textquotedblleft{}an interval\textquotedblright{}; MW gives \emph{nir-antaram} \textquotedblleft{}closely\textquotedblright{} and \textquotedblleft{}constantly, continually\textquotedblright{} (MW).
\end{cousinweb}

\begin{grammarlens}[title={What you've built}]
One more word for saying what happens if.
\end{grammarlens}

\subsection*{Your turn}
\begin{itemize}
\raggedright
  \item \emph{Say it:} \emph{nirantaram}
  \item \emph{Say it:} \emph{nirantaram}, once more
  \item \emph{Say it:} \emph{yadi bhavān nirantaraṃ paṭhati, tarhi parīkṣāyām uttamaḥ bhaviṣyati.} — If you study without a break, you will do well in the exam.
  \item \emph{Recall:} say the Sanskrit for triumphs, then the Sanskrit for unshakable, then say \emph{nirantaram} again
\end{itemize}

\begin{tcolorbox}[breakable,colback=teal!4,colframe=teal!35!black,title={Before you move on}]
What does \sk{निरन्तरम्} mean? (\textquotedblleft{}Without a break; constantly\textquotedblright{}.) Say it once more.
\end{tcolorbox}
\section[\texorpdfstring{\sk{तत्परः} (tatparaḥ)}{तत्परः (tatparaḥ)}]{\sk{तत्परः} (tatparaḥ) — intent on; fully engaged}

\label{lesson:SA-C393-tatparah}



\begin{quote}
Before the new one: say the Sanskrit for unshakable, then the Sanskrit for without a break.
\end{quote}

\subsection*{You'll want to know: \sk{तत्परः}}
\textbf{\sk{तत्परः}} — \emph{tatparaḥ} — \textquotedblleft{}intent on; fully engaged\textquotedblright{}.

\begin{quote}
\sk{यदि} \sk{छात्रः} \sk{अध्ययने} \sk{तत्परः}, \sk{तर्हि} \sk{सः} \sk{सफलतां} \sk{प्राप्नोति।}
\end{quote}

\begin{quote}
\emph{yadi chātraḥ adhyayane tatparaḥ, tarhi saḥ saphalatāṃ prāpnoti.}
\end{quote}

\begin{quote}
If a student is intent on study, he succeeds.
\end{quote}

\begin{cousinweb}
\textbf{\sk{तत्परः}} (\emph{tatparaḥ}) is built from \emph{tat}, \textquotedblleft{}that\textquotedblright{}, and \emph{para}, \textquotedblleft{}highest\textquotedblright{}; MW gives \emph{tat-para} \textquotedblleft{}having that as one's highest object, totally devoted to, eagerly engaged in\textquotedblright{} (MW).
\end{cousinweb}

\begin{grammarlens}[title={What you've built}]
One more word for saying what happens if.
\end{grammarlens}

\subsection*{Your turn}
\begin{itemize}
\raggedright
  \item \emph{Say it:} \emph{tatparaḥ}
  \item \emph{Say it:} \emph{tatparaḥ}, once more
  \item \emph{Say it:} \emph{yadi chātraḥ adhyayane tatparaḥ, tarhi saḥ saphalatāṃ prāpnoti.} — If a student is intent on study, he succeeds.
  \item \emph{Recall:} say the Sanskrit for unshakable, then the Sanskrit for without a break, then say \emph{tatparaḥ} again
\end{itemize}

\begin{tcolorbox}[breakable,colback=teal!4,colframe=teal!35!black,title={Before you move on}]
What does \sk{तत्परः} mean? (\textquotedblleft{}Intent on; fully engaged\textquotedblright{}.) Say it once more.
\end{tcolorbox}
\section[\texorpdfstring{\sk{कृतकृत्यः} (kṛtakṛtyaḥ)}{कृतकृत्यः (kṛtakṛtyaḥ)}]{\sk{कृतकृत्यः} (kṛtakṛtyaḥ) — fulfilled; content with what one has done}

\label{lesson:SA-C393-krtakrtyah}



\begin{quote}
Before the new one: say the Sanskrit for without a break, then the Sanskrit for intent on.
\end{quote}

\subsection*{You'll want to know: \sk{कृतकृत्यः}}
\textbf{\sk{कृतकृत्यः}} — \emph{kṛtakṛtyaḥ} — \textquotedblleft{}fulfilled; content with what one has done\textquotedblright{}.

\begin{quote}
\sk{यदि} \sk{सः} \sk{लक्ष्यं} \sk{प्राप्नोति}, \sk{तर्हि} \sk{कृतकृत्यः} \sk{भविष्यति।}
\end{quote}

\begin{quote}
\emph{yadi saḥ lakṣyaṃ prāpnoti, tarhi kṛtakṛtyaḥ bhaviṣyati.}
\end{quote}

\begin{quote}
If he reaches his goal, he will feel fulfilled.
\end{quote}

\begin{cousinweb}
\textbf{\sk{कृतकृत्यः}} (\emph{kṛtakṛtyaḥ}) is built from \emph{kṛta}, \textquotedblleft{}done\textquotedblright{}, and \emph{kṛtya}, \textquotedblleft{}what is to be done\textquotedblright{}, both from the root √kṛ; MW gives \emph{kṛta-kṛtya} \textquotedblleft{}one who has done his duty\textquotedblright{} and \textquotedblleft{}contented, satisfied\textquotedblright{} (MW).
\end{cousinweb}

\begin{grammarlens}[title={What you've built}]
One more word for saying what happens if.
\end{grammarlens}

\subsection*{Your turn}
\begin{itemize}
\raggedright
  \item \emph{Say it:} \emph{kṛtakṛtyaḥ}
  \item \emph{Say it:} \emph{kṛtakṛtyaḥ}, once more
  \item \emph{Say it:} \emph{yadi saḥ lakṣyaṃ prāpnoti, tarhi kṛtakṛtyaḥ bhaviṣyati.} — If he reaches his goal, he will feel fulfilled.
  \item \emph{Recall:} say the Sanskrit for without a break, then the Sanskrit for intent on, then say \emph{kṛtakṛtyaḥ} again
\end{itemize}

\begin{tcolorbox}[breakable,colback=teal!4,colframe=teal!35!black,title={Before you move on}]
What does \sk{कृतकृत्यः} mean? (\textquotedblleft{}Fulfilled; content with what one has done\textquotedblright{}.) Say it once more.
\end{tcolorbox}
\section[(dialogue)]{Looking back: from Getting Better to What It Takes}

\label{lesson:SA-R393-review-387-393}



\begin{quote}
Say the Sanskrit for advancement and development.
\end{quote}

\subsection*{Your turn}
Say these aloud:

\begin{itemize}
\raggedright
  \item advancement, development, hard times, refines, growth
  \item deals with people, neglects, mild-mannered, in vain, an ideal
  \item possible, achievable, imagines, imagination, true to fact
  \item a sign of something to come, an expectation, useless, worthwhile, experiences
  \item supreme, inferior, in keeping with, suitable, spoils
  \item requests, an assistant, reassures, co-operation, a kindness
  \item triumphs, unshakable, without a break, intent on, fulfilled
\end{itemize}

\begin{tcolorbox}[breakable,colback=teal!4,colframe=teal!35!black,title={Before you move on}]
Advancement? (\textbf{\sk{उन्नतिः}}, \emph{unnatiḥ}.) Development? (\textbf{\sk{विकासः}}, \emph{vikāsaḥ}.) Hard times? (\textbf{\sk{दुर्गतिः}}, \emph{durgatiḥ}.) Refines? (\textbf{\sk{संस्करोति}}, \emph{saṃskaroti}.) Growth? (\textbf{\sk{वृद्धिः}}, \emph{vṛddhiḥ}.) Deals with people? (\textbf{\sk{व्यवहरति}}, \emph{vyavaharati}.) Neglects? (\textbf{\sk{उपेक्षते}}, \emph{upekṣate}.) Mild-mannered? (\textbf{\sk{सौम्यः}}, \emph{saumyaḥ}.) In vain? (\textbf{\sk{वृथा}}, \emph{vṛthā}.) An ideal? (\textbf{\sk{आदर्शः}}, \emph{ādarśaḥ}.) Possible? (\textbf{\sk{शक्यः}}, \emph{śakyaḥ}.) Achievable? (\textbf{\sk{साध्यः}}, \emph{sādhyaḥ}.) Imagines? (\textbf{\sk{कल्पयति}}, \emph{kalpayati}.) Imagination? (\textbf{\sk{कल्पना}}, \emph{kalpanā}.) True to fact? (\textbf{\sk{यथार्थः}}, \emph{yathārthaḥ}.) A sign of something to come? (\textbf{\sk{लक्षणम्}}, \emph{lakṣaṇam}.) An expectation? (\textbf{\sk{प्रत्याशा}}, \emph{pratyāśā}.) Useless? (\textbf{\sk{व्यर्थः}}, \emph{vyarthaḥ}.) Worthwhile? (\textbf{\sk{सार्थकः}}, \emph{sārthakaḥ}.) Experiences? (\textbf{\sk{अनुभवति}}, \emph{anubhavati}.) Supreme? (\textbf{\sk{परमः}}, \emph{paramaḥ}.) Inferior? (\textbf{\sk{निकृष्टः}}, \emph{nikṛṣṭaḥ}.) In keeping with? (\textbf{\sk{अनुरूपः}}, \emph{anurūpaḥ}.) Suitable? (\textbf{\sk{उपयुक्तः}}, \emph{upayuktaḥ}.) Spoils? (\textbf{\sk{दूषयति}}, \emph{dūṣayati}.) Requests? (\textbf{\sk{अभ्यर्थयते}}, \emph{abhyarthayate}.) An assistant? (\textbf{\sk{सहायकः}}, \emph{sahāyakaḥ}.) Reassures? (\textbf{\sk{आश्वासयति}}, \emph{āśvāsayati}.) Co-operation? (\textbf{\sk{सहकारः}}, \emph{sahakāraḥ}.) A kindness? (\textbf{\sk{अनुग्रहः}}, \emph{anugrahaḥ}.) Triumphs? (\textbf{\sk{विजयते}}, \emph{vijayate}.) Unshakable? (\textbf{\sk{अविचलः}}, \emph{avicalaḥ}.) Without a break? (\textbf{\sk{निरन्तरम्}}, \emph{nirantaram}.) Intent on? (\textbf{\sk{तत्परः}}, \emph{tatparaḥ}.) Fulfilled? (\textbf{\sk{कृतकृत्यः}}, \emph{kṛtakṛtyaḥ}.)
\end{tcolorbox}
\section[(dialogue)]{What It Takes — once more}

\label{lesson:SA-R393-review-final-393}



\begin{quote}
Say the Sanskrit for triumphs and unshakable.
\end{quote}

\subsection*{Your turn}
\begin{itemize}
\raggedright
  \item \emph{Say it:} triumphs, unshakable, without a break, intent on, fulfilled
\end{itemize}

\begin{tcolorbox}[breakable,colback=teal!4,colframe=teal!35!black,title={Before you move on}]
Triumphs? (\textbf{\sk{विजयते}}, \emph{vijayate}.) Unshakable? (\textbf{\sk{अविचलः}}, \emph{avicalaḥ}.) Without a break? (\textbf{\sk{निरन्तरम्}}, \emph{nirantaram}.) Intent on? (\textbf{\sk{तत्परः}}, \emph{tatparaḥ}.) Fulfilled? (\textbf{\sk{कृतकृत्यः}}, \emph{kṛtakṛtyaḥ}.)
\end{tcolorbox}
